% Part: first-order-logic
% Chapter: sequent-calculus
% Section: rules-and-proofs

\documentclass[../../../include/open-logic-section]{subfiles}

\begin{document}

\iftag{FOL}
      {\olfileid[fr]{fol}{seq}{rul}}
      {\olfileid[fr]{pl}{seq}{rul}}

\olsection{Règles et \usetoken{p}{derivation}}

Dans ce qui suit, $\Gamma, \Delta, \Pi, \Lambda$ désignent des suites
finies d'!!{sentence}s.

\begin{defn}[Séquent]
Un \emph{séquent} est une expression de la forme
\[
\Gamma \Sequent \Delta
\]
où $\Gamma$ et $\Delta$ sont des suites finies, éventuellement vides,
d'!!{sentence}s du langage $\Lang L$. On appelle $\Gamma$
l'\emph{antécédent} et $\Delta$ le \emph{conséquent}.
\end{defn}

\begin{explain}
L'idée intuitive d'un séquent est la suivante : si tous les
!!{sentence}s de l'antécédent sont vrais, alors au moins un des
!!{sentence}s du conséquent est vrai. Autrement dit, si
$\Gamma = \tuple{!A_1, \dots, !A_m}$ et
$\Delta = \tuple{!B_1, \dots, !B_n}$, alors
$\Gamma \Sequent \Delta$ est vrai si et seulement si
\[
(!A_1 \land \cdots \land !A_m) \lif (!B_1 \lor \cdots \lor
!B_n)
\]
est vrai. Deux cas particuliers se présentent : celui où $\Gamma$
est vide et celui où $\Delta$~est vide. Lorsque $\Gamma$ est vide,
c'est-à-dire lorsque $m = 0$, $\quad \Sequent \Delta$ est vrai si et
seulement si $!B_1 \lor \dots \lor !B_n$ est vrai. Lorsque $\Delta$
est vide, c'est-à-dire lorsque $n = 0$, $\Gamma \Sequent \quad$ est
vrai si et seulement si $\lnot(!A_1 \land \dots \land !A_m)$ est vrai.
On dit qu'un séquent est valide si et seulement si l'!!{sentence}
correspondant est valide.
\end{explain}

Si $\Gamma$ est une suite d'!!{sentence}s, on note $\Gamma, !A$
la suite obtenue en ajoutant $!A$ à l'extrémité droite de~$\Gamma$
(et $!A, \Gamma$ celle obtenue en ajoutant $!A$ à son extrémité
gauche). Si $\Delta$ est aussi une suite d'!!{sentence}s, alors
$\Gamma, \Delta$ désigne la concaténation des deux suites.

\Needspace{9\baselineskip}
\begin{defn}[Séquent initial]
Soit $!A$ un !!{sentence} quelconque du langage.
Un \emph{séquent initial} est un séquent
\iftag{prvFalse,prvTrue}{de l'une des formes suivantes :
  \begin{enumerate}
    \item $!A \Sequent !A$.
    \tagitem{prvTrue}{$\quad \Sequent \ltrue$.}{}
    \tagitem{prvFalse}{$\lfalse \Sequent \quad$.}{}
  \end{enumerate}}
{de la forme $!A \Sequent !A$.}
\end{defn}

Les !!{derivation}s du calcul des séquents sont certains arbres de
séquents. Les séquents situés tout en haut sont des séquents initiaux ;
un séquent placé sous un ou deux autres doit s'en déduire par une
application correcte d'une règle d'inférence. Les règles de
$\Log{LK}$ se répartissent en deux grandes catégories : les règles
\emph{logiques} et les règles \emph{structurelles}. Les règles
logiques portent le nom de l'!!{main operator} de l'!!{sentence} qui
contient $!A$ et/ou $!B$ dans le séquent inférieur. Chacune existe
en deux versions : l'une permet d'obtenir un séquent où
l'!!{sentence} contenant le !!{operator} figure à gauche, l'autre
un séquent où cet !!{sentence} figure à droite.

\end{document}
